\documentclass[12pt,a4paper,titlepage]{article}

\usepackage{float}
\usepackage{xcolor}
\usepackage{amsmath}
\usepackage{titlesec}
\usepackage{graphicx}
\usepackage[portuguese]{babel}
\usepackage[colorlinks=true,allcolors=black]{hyperref}

\graphicspath{ {./img/}{./../../../../} }

\newcommand{\sectionbreak}{\clearpage}

\title{Arquitetura de Computadores II\\Lista de Estudos -- Prova II}
\author{
    \includegraphics[scale=0.30]{puc}
    \\[1.5em]
    \\Fernando de Lima Gonçalves
}
\date{Belo Horizonte\\2026}

\begin{document}
\maketitle
\tableofcontents

\section{Parte C - Múltipla escolha}
\begin{table}[h]
    \centering
    \caption{Gabarito Múltipla escolha}
    \begin{tabular}{|c|c|}
        \hline
        \textbf{Questão} & \textbf{Resposta} \\
        \hline
        01 & D \\
        \hline
        02 & C \\
        \hline
        03 & A \\
        \hline
        04 & B \\
        \hline
        05 & A \\
        \hline
        06 & D \\
        \hline
        07 & C \\
        \hline
        08 & B \\
        \hline
        09 & B \\
        \hline
        10 & D \\
        \hline
        11 & A \\
        \hline
        12 & C \\
        \hline
        13 & D \\
        \hline
        14 & B \\
        \hline
        15 & A \\
        \hline
    \end{tabular}
    \hspace{1cm}
    \begin{tabular}{|c|c|}
        \hline
        \textbf{Questão} & \textbf{Resposta} \\
        \hline
        16 & C \\
        \hline
        17 & A \\
        \hline
        18 & C \\
        \hline
        19 & D \\
        \hline
        20 & B \\
        \hline
        21 & C \\
        \hline
        22 & B \\
        \hline
        23 & B \\
        \hline
        24 & D \\
        \hline
        25 & A \\
        \hline
        26 & C \\
        \hline
        27 & A \\
        \hline
        28 & B \\
        \hline
        29 & A \\
        \hline
        30 & D \\
        \hline
    \end{tabular}
    \label{gabarito_multipla_escolha}
\end{table}

\section{Parte D - Verdadeiro ou Falso}
\begin{table}[h]
    \centering
    \caption{Gabarito Verdadeiro ou Falso}
    \begin{tabular}{|c|c|}
        \hline
        \textbf{Questão} & \textbf{Resposta} \\
        \hline
        01 & 0 \\
        \hline
        02 & 1 \\
        \hline
        03 & 1 \\
        \hline
        04 & 0 \\
        \hline
        05 & 1 \\
        \hline
        06 & 1 \\
        \hline
        07 & 0 \\
        \hline
        08 & 0 \\
        \hline
        09 & 1 \\
        \hline
        10 & 1 \\
        \hline
        11 & 1 \\
        \hline
        12 & 1 \\
        \hline
        13 & 0 \\
        \hline
        14 & 1 \\
        \hline
        15 & 0 \\
        \hline
        16 & 1 \\
        \hline
        17 & 1 \\
        \hline
        18 & 1 \\
        \hline
        19 & 0 \\
        \hline
        20 & 1 \\
        \hline
        21 & 0 \\
        \hline
        22 & 1 \\
        \hline
        23 & 0 \\
        \hline
        24 & 1 \\
        \hline
        25 & 0 \\
        \hline
        26 & 1 \\
        \hline
        27 & 0 \\
        \hline
    \end{tabular}
    \hspace{1cm}
    \begin{tabular}{|c|c|}
        \hline
        \textbf{Questão} & \textbf{Resposta} \\
        \hline
        28 & 1 \\
        \hline
        29 & 1 \\
        \hline
        30 & 0 \\
        \hline
        31 & 0 \\
        \hline
        32 & 1 \\
        \hline
        33 & 1 \\
        \hline
        34 & 1 \\
        \hline
        35 & 0 \\
        \hline
        36 & 1 \\
        \hline
        37 & 1 \\
        \hline
        38 & 1 \\
        \hline
        39 & 0 \\
        \hline
        40 & 1 \\
        \hline
        41 & 0 \\
        \hline
        42 & 0 \\
        \hline
        43 & 1 \\
        \hline
        44 & 0 \\
        \hline
        45 & 1 \\
        \hline
        46 & 0 \\
        \hline
        47 & 0 \\
        \hline
        48 & 1 \\
        \hline
        49 & 1 \\
        \hline
        50 & 0 \\
        \hline
        51 & 1 \\
        \hline
        52 & 1 \\
        \hline
        53 & 0 \\
        \hline
        54 & 0 \\
        \hline
    \end{tabular}
    \label{gabarito_v_f}
\end{table}

Considerando questão 01 como MSB e questão 54 como LSB temos:
\begin{center}
    \fbox{\textbf{Binário:} 011011001111010111010101010110011101110100101001101100}

    \vspace{0.5em}

    \fbox{\textbf{Decimal:} 7.667.398.542.314.092}
\end{center}

\section{Parte E - Estudo Dirigido}
\subsection{ISA e Assembly}
\begin{itemize}
    \item\textbf{E1:}
        A manutenção dos campos rs1, rs2 e rd nas mesmas posições, nos formatos em que estão presentes, permite que a leitura do banco
        de registradores ocorra simultaneamente à decodificação da instrução, sem precisar aguardar a identificação do opcode.
        Isso reduz a necessidade de multiplexadores e simplifica as conexões do caminho de dados.
        Além disso, manter o bit de sinal do imediato na posição 31 permite iniciar a extensão de sinal antes da identificação completa
        do formato da instrução, reduzindo o caminho crítico do monociclo e simplificando o estágio ID do pipeline.

    \item\textbf{E2:}
        Os campos imm[10:5] e imm[4:1] do formato B ocupam as mesmas posições físicas dos respectivos campos do formato S.
        No formato B, imm[12] corresponde ao bit de sinal da instrução, imm[11] ocupa a posição 7 e imm[0] é implicitamente zero.
        Essa organização permite reutilizar grande parte das conexões físicas do formato S, evitando um deslocamento adicional do
        imediato e reduzindo a necessidade de multiplexadores. Embora a codificação seja menos intuitiva para interpretação humana,
        ela simplifica o gerador de imediato e reduz o custo de implementação do hardware.

    \item\textbf{E3:}
        O prólogo utiliza addi para reservar espaço na pilha e sw para salvar os registradores ra e s utilizados pela rotina.
        Sem esse salvamento, a instrução jal que chama a função secundária sobrescreve ra com um novo endereço de retorno.
        Ao executar ret, a função secundária retorna corretamente à rotina principal, mas o ret da rotina principal utiliza o
        endereço sobrescrito, podendo provocar um laço infinito ou uma falha na execução.
        Os registradores s são classificados como callee-saved e devem ser preservados pelas funções que os modificam.
        Caso contrário, valores utilizados como índices ou acumuladores podem ser perdidos. O epílogo utiliza lw para restaurar
        os registradores e addi para recuperar o valor original de sp, respeitando o alinhamento da pilha.

    \item\textbf{E4:}
        Na transposição realizada no próprio espaço de memória, basta percorrer os elementos acima ou abaixo da diagonal principal,
        considerando, por exemplo, a condição $i < j$. Para uma matriz quadrada de ordem $n$, são necessárias
        $\frac{n(n-1)}{2}$ trocas, enquanto a transposição para uma nova matriz exige $n^2$ cópias e espaço adicional em memória.
        Para uma matriz $3 \times 3$, a transposição no lugar exige três trocas, seis instruções lw e seis instruções sw,
        enquanto a transposição para uma nova matriz exige nove instruções lw e nove instruções sw.
        A primeira abordagem utiliza registradores temporários para realizar as trocas, enquanto a segunda necessita manter
        os endereços das duas matrizes. A redução dos acessos à memória e das iterações tende a reduzir o número total de instruções.
\end{itemize}

\subsection{Microarquitetura monociclo}
\begin{itemize}
    \item\textbf{E5:}
        Uma instrução do tipo R, como add, não utiliza a memória de dados nem o imediato. Uma instrução lw não utiliza o segundo
        operando do banco de registradores nem a lógica de desvio. Uma instrução sw não utiliza a escrita no banco de registradores,
        enquanto beq não utiliza a memória de dados nem a escrita no banco. Por fim, jal não utiliza a memória de dados nem os
        operandos lidos do banco de registradores.
        No monociclo, o período de clock é determinado pela instrução de maior caminho crítico, normalmente lw, fazendo com que
        instruções mais simples aguardem um período maior que o necessário. O pipeline reduz essa ineficiência ao permitir que
        diferentes instruções utilizem simultaneamente diferentes blocos do caminho de dados.

    \item\textbf{E6:}
        Na instrução add t0, t0, t1, a leitura combinacional do banco de registradores disponibiliza os valores atuais de t0 e t1.
        A ALU realiza a soma durante o ciclo, mas o novo valor de t0 somente é armazenado na borda ativa do clock, garantindo
        que o cálculo utilize seu valor anterior.
        Em Verilog, esse comportamento é representado pela atribuição não bloqueante, indicada pelo operador \texttt{<=},
        nos blocos sequenciais. No pipeline, também é possível realizar a escrita no banco na primeira metade do ciclo e a
        leitura na segunda metade, permitindo que uma instrução em ID obtenha o valor escrito por outra instrução em WB.

    \item\textbf{E7:}
        O valor X em uma tabela de sinais de controle representa uma condição don't care, indicando que determinado sinal pode
        assumir 0 ou 1 sem modificar o resultado da instrução.
        Em uma instrução sw, por exemplo, RegWrite permanece em 0, tornando irrelevante o valor de MemtoReg.
        Analogamente, ALUSrc é irrelevante quando o resultado produzido pela ALU não é utilizado, dependendo da organização
        do caminho de dados. Entretanto, em instruções como lw e sw, seu valor deve ser definido corretamente para calcular
        o endereço de memória.
        Essas condições permitem simplificações durante a síntese do circuito, mas podem provocar erros caso outros sinais
        estejam incorretamente habilitados. Em Verilog, o valor X também pode representar um estado lógico desconhecido
        durante a simulação.
\end{itemize}

\subsection{Verilog}
\begin{itemize}
    \item\textbf{E8:}
        Para representar o paralelismo do hardware em uma linguagem sequencial, é necessário separar a avaliação da lógica
        combinacional da atualização dos elementos sequenciais. Durante cada ciclo, os resultados intermediários são calculados
        utilizando o estado atual do processador e armazenados temporariamente.
        Somente ao final do ciclo são atualizados o PC, o banco de registradores e a memória de dados, quando habilitados.
        Dessa forma, a ordem das atribuições deve respeitar as dependências combinacionais, sem permitir que uma atualização
        prematura altere os valores utilizados no mesmo ciclo. Esse comportamento corresponde ao das atribuições não bloqueantes
        utilizadas nos blocos sequenciais do Verilog.

    \item\textbf{E9:}
        Uma sequência de valores do PC igual a 0, 4, 8 e 0 pode representar tanto um laço válido quanto um erro de implementação.
        Deve-se verificar o conteúdo da ROM de instruções, principalmente a instrução armazenada no endereço 8, identificando
        se existe uma instrução jal ou um desvio condicional tomado cujo destino seja o endereço 0.
        Nesse caso, um desvio relativo ao PC deve possuir deslocamento igual a -8.
        Também devem ser analisados os sinais Branch, Jump e Zero, o imediato gerado, o reset e a seleção do próximo PC.
        Testes isolados de incremento, reset e desvio permitem comparar os resultados obtidos com os valores esperados.

    \item\textbf{E10:}
        Um testbench que utiliza exclusivamente o comando \texttt{\$display} apresenta os resultados produzidos pelo módulo,
        mas não verifica automaticamente sua corretude. Para testar a unidade de controle, deve-se utilizar uma tabela-verdade
        que relacione cada opcode aos respectivos sinais de controle esperados.
        O testbench deve percorrer os opcodes suportados, comparar os sinais produzidos com os valores esperados, desconsiderando
        aqueles classificados como don't care, e contabilizar as divergências em um contador de erros.
        Também devem ser testados opcodes inválidos ou não implementados, verificando se operações indevidas de escrita
        são desabilitadas. Essa abordagem caracteriza um self-checking testbench.
\end{itemize}

\clearpage\subsection{Pipeline}
\begin{itemize}
    \item\textbf{E11:}
        Para transformar um processador monociclo em um pipeline de cinco estágios, é necessário dividir o caminho de dados
        nas etapas IF, ID, EX, MEM e WB, inserindo os registradores intermediários IF/ID, ID/EX, EX/MEM e MEM/WB.
        Esses registradores armazenam os operandos, os resultados intermediários, o registrador de destino rd e os sinais
        de controle de cada instrução. Como diversas instruções são executadas simultaneamente, os sinais de controle precisam
        acompanhar suas respectivas instruções para serem utilizados nos estágios corretos, como MemWrite em MEM e RegWrite em WB.
        Também devem ser acrescentados mecanismos de forwarding, detecção de conflitos, inserção de bolhas e descarte de instruções.

    \item\textbf{E12:}
        A sequência formada pelas instruções lw t2, 0(t0) e add a0, a0, t2 apresenta um conflito de dados do tipo load-use hazard,
        pois o valor carregado por lw somente estará disponível ao final do estágio MEM, enquanto add necessita desse valor
        no estágio EX. Mesmo com forwarding, é necessária a inserção de uma bolha no pipeline clássico de cinco estágios.

        \begin{center}
            \textbf{Sem reordenação}

            \smallskip
            \begin{tabular}{|l|*{8}{c|}}
                \hline
                Instrução & 1 & 2 & 3 & 4 & 5 & 6 & 7 & 8 \\
                \hline
                lw t2, 0(t0)    & IF & ID & EX & MEM   & WB &     &     &    \\
                add a0, a0, t2  &    & IF & ID & Stall & EX & MEM & WB  &    \\
                addi t0, t0, 4  &    &    & IF & Stall & ID & EX  & MEM & WB \\
                \hline
            \end{tabular}
        \end{center}

        O compilador pode posicionar uma instrução independente, como addi t0, t0, 4, entre lw e add, ocupando o ciclo
        anteriormente desperdiçado.

        \begin{center}
            \textbf{Com reordenação}

            \smallskip
            \begin{tabular}{|l|*{8}{c|}}
                \hline
                Instrução & 1 & 2 & 3 & 4 & 5 & 6 & 7 & 8 \\
                \hline
                lw t2, 0(t0)    & IF & ID & EX & MEM & WB  &     &    & \\
                addi t0, t0, 4  &    & IF & ID & EX  & MEM & WB  &    & \\
                add a0, a0, t2  &    &    & IF & ID  & EX  & MEM & WB & \\
                \hline
            \end{tabular}
        \end{center}

        Essa reorganização reduz o CPI sem modificar a quantidade de instruções ou a semântica do programa.
        Entretanto, deve-se respeitar as dependências de dados, evitando modificar t0 antes da leitura realizada por lw,
        pois isso alteraria o endereço acessado na memória.

    \item\textbf{E13:}
        O aumento da quantidade de estágios do pipeline permite dividir o caminho crítico em etapas menores, possibilitando
        a redução do período de clock e o aumento da frequência de operação.
        Entretanto, pipelines mais profundos podem apresentar penalidades maiores em erros de predição de desvios, pois uma
        quantidade maior de instruções especulativas pode precisar ser descartada.
        Considerando a equação $T = \frac{I_c \times CPI}{f}$, o aumento da frequência somente melhora o desempenho quando
        compensa o eventual aumento do CPI. Para um mesmo número de instruções, a razão $CPI/f$ determina o tempo de execução.
        Eventos como branch-misses, disponibilizados pelo perf, permitem investigar a ocorrência de erros de predição.
\end{itemize}

\subsection{Avaliação de desempenho}
\begin{itemize}
    \item\textbf{E14:}
        O código C é compilado diretamente para instruções de máquina, enquanto Python, considerando a implementação CPython,
        executa bytecode por meio de um interpretador. Uma única operação em Python pode exigir diversas instruções de máquina,
        aumentando significativamente $I_c$.
        Java também apresenta custos associados à inicialização da JVM e à interpretação inicial do bytecode. Além disso,
        diferenças nos acessos à memória e na execução de desvios podem influenciar o CPI.
        Caso a busca seja repetida milhares de vezes no mesmo processo Java, o compilador JIT poderá converter o método
        frequentemente executado em código de máquina otimizado, reduzindo $I_c$ e podendo melhorar o CPI.
        Esse processo de aquecimento, denominado warm-up, pode aproximar o desempenho de Java ao de C, reduzindo a diferença
        observada no experimento.

    \item\textbf{E15:}
        Quando o perf realiza a medição do programa inteiro, são contabilizadas tanto as instruções do algoritmo de busca
        quanto aquelas relacionadas à leitura de arquivos, alocação de memória e inicialização do programa.
        Como as instruções da busca representam menos de 1\% do total, as diferenças entre o melhor e o pior caso podem ser
        mascaradas pelas demais operações e pelo ruído do sistema.
        O grupo poderia concluir incorretamente que os dois casos apresentam desempenhos semelhantes, apesar de possuírem
        complexidades distintas, respectivamente $\Theta(1)$ e $\Theta(n)$.
        Para melhorar o experimento, deve-se carregar previamente os dados em memória, repetir a busca diversas vezes e
        restringir a medição à região de interesse, utilizando ferramentas como Cachegrind e perf record para investigar
        os eventos associados às funções.

    \item\textbf{E16:}
        Mesmo que dois compiladores produzam laços com a mesma quantidade de instruções por iteração, seus tempos de execução
        podem apresentar diferenças decorrentes de características microarquiteturais.
        O alinhamento das instruções pode influenciar a utilização da cache de instruções, enquanto o sentido e a organização
        dos desvios podem alterar a eficiência da predição.
        Além disso, determinados processadores realizam a macrofusão, combinando instruções de comparação e desvio em uma
        única operação interna, sem modificar a quantidade de instruções arquiteturais.
        Dependências de dados, variações na frequência de operação e interferências do sistema operacional também podem
        modificar o tempo observado. Portanto, a igualdade de $I_c$ não implica necessariamente igualdade de CPI.

    \item\textbf{E17:}
        A média aritmética de tempos normalizados depende da máquina utilizada como referência, pois a normalização modifica
        o peso relativo de cada benchmark. Considere duas máquinas A e B, com tempos de execução de 1 e 3 segundos para A
        e 2 e 1 segundos para B.
        Normalizando pela máquina A, as médias aritméticas são 1 e aproximadamente 1,167, favorecendo A.
        Normalizando pela máquina B, as médias passam a ser 1,75 e 1, favorecendo B.
        A média geométrica não apresenta esse problema, pois satisfaz a relação
        $\frac{MG(A/R)}{MG(B/R)} = \frac{MG(A)}{MG(B)}$.
        Dessa forma, o termo referente à máquina de referência é cancelado, mantendo a classificação relativa das máquinas
        independentemente da normalização escolhida.
\end{itemize}

\subsection{Microcontroladores}
\begin{itemize}
    \item\textbf{E18:}
        A arquitetura Harvard, utilizada pelos microcontroladores AVR, caracteriza-se pela separação das memórias e dos barramentos
        destinados às instruções e aos dados. Essa organização permite realizar a busca da próxima instrução simultaneamente
        à execução da instrução atual, caracterizando o pipeline descrito no datasheet.
        Em um pipeline RISC-V de cinco estágios que utilizasse uma memória única com apenas uma porta de acesso, os estágios IF
        e MEM poderiam disputar simultaneamente o mesmo recurso, caracterizando um conflito estrutural.
        A arquitetura Harvard evita esse conflito ao disponibilizar caminhos independentes de acesso, de maneira equivalente
        à separação das caches L1 de instruções e de dados em processadores modernos.

    \item\textbf{E19:}
        O consumo energético de um microcontrolador depende tanto da potência consumida quanto do tempo necessário para realizar
        determinada atividade. O tempo ativo é dado por $T_{\mathrm{ativo}} = \frac{I_c \times CPI}{f}$.
        A energia consumida durante um período de medição pode ser estimada por
        $E = P_{\mathrm{ativo}} \times T_{\mathrm{ativo}} + P_{\mathrm{sono}} \times T_{\mathrm{sono}}$.
        Um ESP32 operando a 240 MHz pode apresentar maior potência durante a execução, mas concluir o processamento mais rapidamente
        e permanecer por mais tempo em modo de baixo consumo, estratégia denominada race to sleep.
        Por outro lado, um ATtiny85 operando a 8 MHz pode apresentar menor potência instantânea, embora necessite de mais tempo
        para executar determinadas tarefas. Também devem ser considerados os custos das transições entre modos de operação
        e da comunicação Wi-Fi. Portanto, a comparação deve considerar a energia total consumida por medição, e não apenas o clock.

    \item\textbf{E20:}
        O SAP-1 possui uma arquitetura didática simplificada, com um conjunto reduzido de instruções, poucos registradores,
        componentes discretos e ausência de pipeline.
        O Intel 8080 é um microprocessador CISC de 8 bits integrado em um único chip, com maior variedade de instruções e
        registradores, mas sem cache integrada ou pipeline nos moldes dos processadores modernos.
        Um núcleo RISC-V atual utiliza uma ISA modular do tipo load/store, podendo incorporar pipelines, caches, predição de
        desvios e execução superescalar, além de integrar outros componentes em um mesmo chip.
        Essas diferenças são justificadas pelos objetivos de cada arquitetura e pela evolução da tecnologia de fabricação,
        da quantidade de transistores disponíveis e da capacidade de memória. Entretanto, pipeline e cache constituem
        decisões microarquiteturais, não exigências da ISA RISC-V.
\end{itemize}

\subsection{Integração entre temas}
\begin{itemize}
    \item\textbf{E21:}
        A instrução lw t2, 8(sp) pertence ao formato I da ISA RV32I e possui código de máquina hexadecimal 0x00812383.
        Seus campos correspondem a imm = 000000001000, rs1 = 00010 (sp), funct3 = 010, rd = 00111 (t2) e opcode = 0000011.
        No estágio IF, o PC fornece o endereço da instrução à memória de instruções. No estágio ID, ocorre a decodificação,
        a leitura de sp e a extensão de sinal do imediato. No estágio EX, a ALU calcula o endereço efetivo mediante a soma
        do conteúdo de sp com o imediato 8. No estágio MEM, ocorre a leitura da palavra armazenada no endereço calculado e,
        no estágio WB, o multiplexador MemtoReg seleciona o dado lido para escrita em t2.
        Os sinais de controle correspondem a RegWrite = 1, ALUSrc = 1, MemRead = 1, MemWrite = 0, MemtoReg = 1, Branch = 0,
        Jump = 0 e ALUOp = 00, considerando uma unidade de controle convencional.
        Em Verilog, os módulos instruction\_memory, register\_file, imm\_gen, control\_unit e ex\_stage são responsáveis
        pela busca, leitura dos registradores, geração do imediato, controle e cálculo do endereço, respectivamente.
        Essas operações podem ser implementadas mediante atribuições assign ou blocos always combinacionais.
        A memória de dados fornece a palavra armazenada no endereço calculado, enquanto a escrita no banco de registradores
        ocorre em um bloco \texttt{always @(posedge clk)}, utilizando uma atribuição não bloqueante e condicionada a RegWrite.

    \item\textbf{E22:}
        Considerando a equação $T = \frac{I_c \times CPI}{f}$, o aumento da frequência de clock reduz proporcionalmente
        o tempo de execução quando a quantidade de instruções e o CPI permanecem constantes.
        Entretanto, frequências maiores podem exigir alterações microarquiteturais, como o aprofundamento do pipeline,
        aumentando as penalidades decorrentes de conflitos e erros de predição.
        Além disso, a memória principal não necessariamente acompanha o aumento da frequência do processador, fazendo com
        que uma mesma latência de memória corresponda a uma quantidade maior de ciclos de clock.
        De acordo com a Lei de Amdahl, a aceleração global é limitada pela parcela do programa efetivamente beneficiada
        pela otimização.
        Portanto, para avaliar o desempenho do programa da Lista V, é necessário considerar $I_c$, CPI, frequência efetiva
        e hierarquia de memória, realizando medições repetidas, com resolução adequada e sob as mesmas condições experimentais.

    \item\textbf{E23:}
        O RISC-V possui uma ISA aberta e modular, permitindo que diferentes fabricantes desenvolvam processadores sem a
        necessidade de pagamento de royalties pelo uso da ISA.
        A base RV32I disponibiliza um conjunto fundamental de instruções inteiras, suficiente para a implementação de
        processadores simples e de baixo custo, podendo ser complementada por extensões como M, F, C e V, conforme
        as necessidades da aplicação.
        Microcontroladores podem utilizar implementações monociclo ou pipelines curtos, priorizando o custo de implementação
        e a eficiência energética. Processadores de alto desempenho podem utilizar pipelines mais complexos, execução
        superescalar, execução fora de ordem, caches e mecanismos sofisticados de predição de desvios.
        Dessa forma, a independência entre ISA e microarquitetura permite que processadores compatíveis com um mesmo
        conjunto de instruções apresentem organizações internas, desempenhos e consumos energéticos distintos.
\end{itemize}

\end{document}
